\documentclass[a4paper]{article}
\usepackage{amsfonts}
\usepackage{amsmath}
\usepackage{amssymb}
\usepackage{graphicx}  
\usepackage{cancel}

\newcommand{\degree}{^{\circ}}
\newcommand{\cis}{\operatorname{cis}}
\newcommand{\Bgsin}{\operatorname{Bgsin}}
\newcommand{\Bgcos}{\operatorname{Bgcos}}
\newcommand{\Bgtan}{\operatorname{Bgtan}}
\newcommand{\limit}{\lim\limits}
\newcommand{\summation}{\displaystyle\sum}
\newcommand{\dom}{\operatorname{dom}}
\newcommand{\Int}{\displaystyle\int}

\begin{document}
  
\noindent \large Auteur: Rune Suy \\
\noindent \large Studierichting: Informatica\\
\noindent \large Oefeningenreeks 5B nr. 5.3\\

\medskip

\normalsize

$\Int_{-1-}^0 \dfrac{dx}{\sqrt{1+x}} = \limit_{a\rightarrow-1-} \Int_{a}^0 \dfrac{dx}{\sqrt{1+x}} = \limit_{a\rightarrow-1-} \left[2\sqrt{1+x}\right]_a^0$\\

$= \limit_{a\rightarrow-1-} 2 - 2\sqrt{1+a} = 2$\\

$\Int \dfrac{dx}{\sqrt{1+x}}$\\

Stel $t=1=x \Rightarrow dt=dx$\\

$= \Int t^{\tfrac{-1}{2}}dt = 2\sqrt{t}+c = 2\sqrt{1+x}+c$\\

\begin{thebibliography}{999}
%\bibitem{a} Referentie
\end{thebibliography}
\end{document} 
